% claim-ref: cherry-shrimp#selected-identity
% claim-ref: cherry-shrimp#adult-size
% claim-ref: cherry-shrimp#preferred-temperature
% claim-ref: cherry-shrimp#broad-temperature
% claim-ref: cherry-shrimp#grazing
% claim-ref: cherry-shrimp#feeding-plan
% claim-ref: cherry-shrimp#molt-cover
% claim-ref: cherry-shrimp#medicine-review
% claim-ref: cherry-shrimp#individual-lifespan
% claim-ref: cherry-shrimp#current-water
\CompanionHeader{ANIMAL CARE}{Cherry shrimp}{Neocaridina davidi / selected red cherry line}{A selected species; stock, age and actual water remain unverified.}
\Context{Daniel is considering a shrimp-only 10-gallon nursery and later adult transfers to the existing 20-gallon Kuhli tank. Neither plan establishes cycling, capacity or compatibility.}
\CardRow{Identity and size}{Red cherry is a color form of N. davidi. UF/IFAS reports maximum adult length around 3--4 cm; females tend to be larger. Confirm the supplier's species and line. \textbf{[CS-UF]}}{Temperature}{Co-Op prefers 72--76 F (22--24 C). Its broad 60--82 F (16--28 C) range is not a setpoint or permission for rapid swings. Source Celsius values are rounded. \textbf{[CS-COOP]}}
\CardRow{Grazing is daily work}{Shrimp browse biofilm and organic material. Keep mature surfaces available; a visually clean, newly filled tank may lack food. Grazing does not replace filtration or planned feeding. \textbf{[CS-UF; CS-COOP]}}{Feed and observe}{Proposed: small varied shrimp-food meals; check access and remove uneaten excess. Offer fine food sparingly near juvenile cover. No universal ration is established. \textbf{[CS-FARM; CS-COOP]}}
\CardRow{Molting and shelter}{Growth requires shedding the shell. Cast shells may be eaten. Soft shrimp need refuge; a molt problem alone does not diagnose a mineral deficiency. Test water and review diet before additives. \textbf{[CS-UF; CS-COOP; proposed]}}{Medicine review}{Review past and proposed treatments. Copper medication can harm invertebrates and biofiltration. Fish-safe does not establish shrimp-safe; seek aquatic veterinary advice for illness. \textbf{[CS-VET; proposed]}}
\Panel{Observe the animals, not just the colony count}{Proposed routine: record grazing, food access, molts, body condition and unusual losses. Check water and equipment if behavior changes. Individual lifespan is unknown here; a continuing colony does not prove how long one shrimp lived.}
\Panel{Before a stocking decision}{Species and line: \hrulefill\quad Supplier water: \hrulefill\par
Health / treatment history: \hrulefill\par
Nursery cycle evidence and mature food surfaces: \hrulefill}
\Panel{Care notes}{Date: \hrulefill\quad Feeding / molts: \hrulefill\par
Observed change / next check: \hrulefill}
\Sources{Sources (clickable): \href{https://ask.ifas.ufl.edu/publication/IN1301}{CS-UF: UF/IFAS species profile}; \href{https://www.aquariumcoop.com/blogs/aquarium/cherry-shrimp-care}{CS-COOP: Co-Op cherry care}; \href{https://www.theshrimpfarm.com/posts/neocaridina-shrimp-care-breeding/}{CS-FARM: Shrimp Farm care}; \href{https://www.msdvetmanual.com/exotic-and-laboratory-animals/aquarium-fish/management-of-aquarium-fish}{CS-VET: MSD veterinary care}. Scopes, derivations and unknowns: adjacent YAML worksheets.}{cherry-shrimp / animal-care}
